\begin{textsample}{POS Dim 4 – rn – Score 6.00 – rn/rn\_em\_13.txt}  \label{ex:f4_pos_059}
3.4 EDUCAÇÃO ESCOLAR \textbf{INDÍGENA} A educação escolar \textbf{indígena} é alicerçada a partir do respeito à etnicidade, à interculturalidade, ao bilinguismo, às diversidades e às especificidades desses povos, e deve ser pensada a partir de um currículo diferenciado e programas específicos advindos de políticas públicas.
Dessa forma, são ratificadas as identidades étnicas, valorizando suas línguas maternas, saberes e tradições, assegurando às comunidades \textbf{indígenas} o direito ao acesso aos conhecimentos técnicos e científicos, a fim da otimização da qualidade de vida, garantindo um espaço para o ensino da língua \textbf{indígena}.
Nesse sentido, reafirmamos a importância das identidades dos povos \textbf{indígenas}, como uma das principais formas de resistência e garantia para manter viva a memória desses povos com cultura, valores, costumes, crenças e saberes, de peculiaridades ímpares.
No Estado do Rio Grande do Norte, registra -se a existência de 13 ( treze ) escolas e creches \textbf{indígenas}, situadas nas comunidades do Amarelão, Santa Terezinha, Serrote, Lagoa do Mato, Tapará, Ladeira Grande e Catu, distribuídas no \textbf{território} Potiguar.
Elencamos, a seguir, algumas definições para melhor compreensão da modalidade : a ) A Etnicidade : compreendida como manifestações de caráter social, uma vez que suas características incluem a língua e a cultura, dentre outras, em que sua transmissão se dá de forma contínua entre os povos de diferentes gerações, configurando -se através do contato e diálogo em sociedade.
Podemos dizer que se refere ao conjunto de características específicas de um grupo de pessoas, com aspectos semelhantes entre si e distintas de outro grupo. b ) A interculturalidade : pode ser compreendida como proposta de convivência democrática, no partilhar de culturas diferentes que, de modo dialógico, convivem sem anular sua diversidade, ressaltando as relações de alteridade ( MONROE, 2009 ).
A proposta pedagógica da escola se faz sob a perspectiva da interculturalidade, buscando assegurar um espaço de reflexão em que os povos \textbf{indígenas} possam dialogar entre as diferentes culturas, resguardando e afirmando as identidades étnicas dos diferentes povos. c ) O bilinguismo e o multilinguismo : respectivamente, convivência de dois sistemas linguísticos diferentes e conhecimento de línguas diversas pelos mesmos falantes, assegurados aos \textbf{indígenas} no texto constitucional de 1988, no art. 231, bem como na Resolução CNE/CEB no 05, de 22 de junho de 2012, que destaca a importância da existência, valorização e manutenção das línguas autóctones, suas crenças, tradições em interação com a língua portuguesa na escola.
Mesmo os povos \textbf{indígenas}, que são hoje monolíngues em língua portuguesa, continuam a usar a língua de seus ancestrais em suas tradições culturais, conhecimentos transmitidos de geração a geração, crenças, práticas religiosas, enfim, as construções socioculturais das sociedades \textbf{indígenas}.
Desse modo, o Referencial Curricular do Ensino Médio Potiguar, em consonância com a Lei de Diretrizes e Bases da Educação Nacional/LDB no 9.394/96 e com a Resolução CNE/CEB no 05, de 22 de junho de 2012, reafirma a importância da reconstrução das memórias históricas dos povos \textbf{indígenas}, suas identidades étnicas, a valorização de suas línguas e ciências, promovendo o desenvolvimento de programas integrados de ensino e pesquisa, possibilitando o acesso às informações, conhecimentos técnicos, científicos e culturais da sociedade nacional e demais sociedades \textbf{indígenas} e não-indígenas, reconhecendo os povos \textbf{indígenas} como sujeitos de direitos.

% matched lemmas: indígena, território
\end{textsample}
